\section{Frobenius orbits as actual local places}
\label{sec:local-identifications}

\status{LI1--LI8 construct the literal local-place identifications,
completion maps, lattice frames and integral away character for supplied
host and level data. The position identity uses the original coefficient
embedding and every prime above $p$. The remaining global comparison
lattice and class-rank theorems are separate inputs.}

\paragraph{The actual arithmetic setting.}
The host has
$G=\operatorname{Gal}(E/\mathbb Q)=(\{\pm1\}^2)\times A$,
$c=((-1,-1),1)$, $F=E^{\langle c\rangle}$ and
$\Psi=\{((\epsilon,1),a)\}$, so $G=\Psi\times\langle c\rangle$.
Write $\sigma_g$ for its action on $E$ and
$\phi_g=\phi_0\sigma_g:E\to\mathbb C$.
Assume $p$ is unramified in $E$ and its arithmetic Frobenius at every
prime above $p$ is
\[
 f=((-1,-1),h)=zc,\qquad z=((1,1),h)\in\Psi,
 \qquad n=\operatorname{ord}(h)\ge2\text{ is even}.       \tag{LI.1}
\]
Put $W_n=W(\mathbb F_{p^n})$, $Q_n=\operatorname{Frac}(W_n)$,
and let $\operatorname{Fr}$ reduce to $x\mapsto x^p$.
LI7 below checks these hypotheses from the literal \texttt{GIdx}
and the supplied host Frobenius law.

\paragraph{LI1--LI2: construct and classify the local embeddings.}
For an unramified prime $P$ of $E$, its completed integer ring has
uniformizer $p$ and residue field $\mathbb F_{p^m}$. Here is the needed
explicit identification with $W(\mathbb F_{p^m})$.
Choose a residue generator $\alpha$, lift its monic minimal polynomial
to $P_0\in\mathbb Z_p[T]$, and lift its simple root by
$b\mapsto b-P_0(b)/P_0'(b)$. The error squares at each step.
Successive residue subtraction and division by $p$ show that
$1,b,\ldots,b^{m-1}$ span the completed ring over $\mathbb Z_p$.
Dividing a hypothetical relation by its smallest $p$-power and reducing
modulo $p$ proves independence. The same construction in the Witt ring
gives the required isomorphism and its uniqueness for a prescribed
residue map. Consequently there is an embedding $i:E\to Q_m$ with
\[
 i(\mathcal O_E)\subset W_m,\qquad
 \overline i:\mathcal O_E\twoheadrightarrow\mathbb F_{p^m},
 \qquad\ker\overline i=P.                               \tag{LI.2}
\]
These are exactly the source's \texttt{Good} conditions.

Two such embeddings inducing $P$ differ by a power of Frobenius:
their residue maps do, and simple-root uniqueness lifts that equality.
Normality of $E/\mathbb Q$ also shows
$\operatorname{Fr}i=i\sigma_f$. Thus $f^m=1$, whereas residue
surjectivity and $f^n=1$ give $\operatorname{Fr}^n=1$ on
$\mathbb F_{p^m}$. The finite-field Frobenius has exact order $m$,
so $m=n$. Fix such an $i$ and put $i_g=i\sigma_g$ and
$P_g=\ker(\overline{i_g})$. We have
\[
 P_{g'}=P_g\ \Longleftrightarrow\ g'=gf^s\text{ for some }s.
                                                               \tag{LI.3}
\]
Every prime above $p$ occurs: construct its embedding by (LI.2),
then use normality to compare it with $i$.

The completion comparison is exact. For any embedding satisfying
(LI.2), $i(x)\in W_n$ if and only if $x$ is integral at its prime:
local denominators outside the prime map to units; for a nonintegral
$x$, its inverse maps to the maximal ideal. Residue lifting approximates
every Witt vector by global integers modulo $p^r$. Global integers are
also dense in the local completion, since a denominator outside the
prime is invertible modulo each prime power. Applying the integrality
criterion to $x/p^r$ identifies the quotients at every precision.
Their inverse limits, and then fraction fields, are therefore isomorphic.

\paragraph{LI3: the orbits label all real-subfield primes.}
The evenness in (LI.1) gives $c\notin\langle f\rangle$: an equality
$f^s=c$ would force both $n\mid s$ and $s$ odd. Hence
$P_g\ne P_{gc}$. Let $v_g=P_g\cap\mathcal O_F$.
Given any residue at $P_g$, choose by CRT an integer $b$ attaining it
and vanishing at $P_{gc}$. Then $b+c(b)\in\mathcal O_F$ has that
same residue. Thus $i_g|_F$ still satisfies (LI.2), and
\[
 Q_n\simeq F_{v_g},\qquad W_n\simeq\mathcal O_{F_{v_g}}. \tag{LI.4}
\]
In particular these primes split in $E/F$.

If $v_{g'}=v_g$ and $y\in P_{g'}$, then $yc(y)\in v_g$,
so $y\in P_g\cup P_{gc}$. An ideal contained in a union of two
ideals lies in one; maximality gives $P_{g'}=P_g$ or $P_{gc}$.
Together with (LI.3), this proves
\[
 \boxed{\ \Psi/\langle z\rangle\ \simeq\
          \{v\subset\mathcal O_F:v\mid p\}\ },
 \qquad\#\{v\mid p\}=|\Psi|/n.                         \tag{LI.5}
\]
Surjectivity uses lying over and the already proved surjectivity for
$E$-primes. Use the literal representatives $a_v=\texttt{orbRep}(z,v)$;
the source positions are $a_vz^s$, $0\le s<n$.

\paragraph{LI4: align the original complex and coefficient embeddings.}
Let $L$ be the coefficient number field, $e_C:L\to\mathbb C$, and
$e:L\to K=\texttt{KfRec}$ the prescribed-place embedding from BC2.
The finite extension $K/\mathbb Q_p$ embeds in
$\overline{\mathbb Q}_p\subset\mathbb C_p$; choose $i_K$.
Integrality of monic roots and their unit inverses shows
$i_K^{-1}(\{|x|\le1\})=\texttt{ORec}$.
Extend $i_Ke$ to an embedding
$\lambda:\overline{\mathbb Q}_{\mathbb C}\to\overline{\mathbb Q}_p$,
so $i_K e=\lambda e_C$ on $L$.
With $\iota:Q_n\to K$, normality of $E$ lets us replace $i$ by
one Galois conjugate so that $i_K\iota i=\lambda\phi_0$.
The preceding construction then uses this adjusted $i$. For
$i_v=i\sigma_{a_v}$ it gives
\[
 i_K\iota\operatorname{Fr}^{s}(i_v(x))
   =\lambda\phi_{a_vz^sc^s}(x).                         \tag{LI.6}
\]
This is the literal \texttt{PosCompat}. The factor $c^s$ is essential:
odd positions conjugate $x\in E$. Every argument of $\lambda$ is
algebraic; no extension to transcendental complex numbers is used.

\paragraph{LI5: the splitting and frame use the given lattice.}
Let $e_v:Q_n\to F_v$ be (LI.4), and $j=e_vi_v$. Choose
$0\ne b\in E$ with $c(b)=-b$, as supplied by the host's
$\sqrt{-d}$. Since $1,b$ is an $F$-basis,
\[
 \Phi_v(x\otimes t)=(j(x)t,j(c(x))t)                    \tag{LI.7}
\]
is an isomorphism $E\otimes_FF_v\simeq F_v\times F_v$:
its inverse sends $(u,v)$ to
$1\otimes(u+v)/2+b\otimes(u-v)/(2j(b))$.
It sends scalars diagonally and conjugation to coordinate interchange.

Write $\Lambda=\texttt{Lc.Lb}\subset E^k$ for the supplied finite
full integer lattice. The first projection of the literal localized
lattice is exactly
\[
 \pi_1(\texttt{locL}_v\Lambda)
   =M_v:=\mathcal O_{F_v}\text{-span}\,j(\Lambda).        \tag{LI.8}
\]
Indeed the source integral tensor subring projects into
$\mathcal O_{F_v}$ by integrality, and contains all its diagonal
scalars. The two inclusions in (LI.8) follow from the defining spans.
The module $M_v$ is finite, torsion free and full: an $E$-basis chosen
from $\Lambda$ retains nonzero determinant under $j$.
Choose a DVR basis as the columns of $B_v\in\mathrm{GL}_k(F_v)$.
Then $t\in M_v$ exactly when $B_v^{-1}t\in\mathcal O_{F_v}^k$.
This proves the exact \texttt{fr\_lat} clause. Equations (LI.4)--(LI.8)
give every field of \texttt{LocId}, with its recovered
$\texttt{iw}(v)=i_v$.

\paragraph{LI6: the same positions induce the prescribed valuations.}
Let $\theta_g:E\to L$ be the actual \texttt{tauLL} embedding with
$e_C\theta_g=\phi_g$. Injectivity of $i_K$ in (LI.6) gives
\[
 e\theta_{a_vz^s}(c^sx)=\iota\operatorname{Fr}^s(i_v(x)).
                                                               \tag{LI.9}
\]
For $0\ne x\in F$, BC3 therefore gives the exact normalization
\[
 \frac{\nu_w(\theta_{a_vz^s}(x))}{\nu_w(p)}=v_p(i_v(x)),
 \qquad
 x\in v\ \Longleftrightarrow\
 \nu_w(\theta_{a_vz^s}(x))>0\quad(x\in\mathcal O_F).
                                                               \tag{LI.10}
\]
The last equivalence includes $x=0$, whose valuation is $+\infty$.
Thus ramified coefficient extension changes the integer valuation scale
by the stated factor; it does not alter which prime is represented.

\paragraph{LI7--LI8: the actual application and integral character.}
For \texttt{GIdx}, the class \texttt{inertClass} excludes the conductor.
The supplied host law gives trivial inertia, hence unramifiedness, and
Frobenius $((-1,-1),h)$ at every prime. Its \texttt{HNeZero} clause
makes $\operatorname{ord}(h)$ even and positive. These are precisely
(LI.1), with $n=\texttt{frobOrder}$. LI1--LI6 therefore prove the actual
\texttt{exists\_locId} conclusion, including position compatibility
for every choice of away character.

Here is that character under the literal \texttt{Joint0} and
\texttt{LevelZeroCharacterRational} assumptions. The local subgroup
equality comes from \texttt{Joint0}; the matching Gram identity is
\texttt{Cmp.gramN\_zero}. For an away-level
element $a$, set its components above $p$ to $1$ and retain all other
components, obtaining $K(a)$ in the original level subgroup $U$.
Only finitely many components change; the unitary equations and level
conditions hold componentwise. Injectivity into the local components
makes $K$ a homomorphism. At level zero, the character
$\chi:U\to\mathbb C^\times$ is a product of additive characters of
finite residue groups. Their values are roots of unity. Rationality
and injectivity of $e_C$ give a multiplicative $r:U\to L^\times$
with $e_Cr=\chi$; both $r(\kappa)$ and its inverse are algebraic
integers. Hence
\[
 \psi(a)=e(r(K(a)))^{-1}\in\texttt{ORec}^{\times},
 \qquad \psi(a)e(r(K(a)))=1.                            \tag{LI.11}
\]
Uniqueness of the componentwise truncation proves this for every
$\kappa$ allowed in \texttt{PsiTied}, establishing \texttt{exists\_psi}.

These choices can be made for all supplied words and places. A common
global frame still uses \texttt{StdTied} and \texttt{Companion}; the
class-rank equality and the actual global comparison-lattice theorem
also remain separate. Neither the existence of the host nor those
global joins follows merely from the local identifications proved here.
